% =====================================================================
\clearpage
\part{Numerics, the missing machine, and the verdict}
% =====================================================================
\section{Numerics and diagnostics}
\label{ch:numerics}
% =====================================================================

\subsection{What the numerics can and cannot show}

The finite-volume corollary of Chapter~\ref{ch:weakstar} sets the frame. At finite $L$, the
transfer matrix is compact and the spectrum is discrete, so $f_1^2>0$ is trivially true. A
finite-lattice measurement of the residue is a statement about a finite-dimensional operator, not
about the continuum theory. The Clay content is the $a\to0$, $L\to\infty$ survival and isolation
of the atom. So the numerics in this chapter are read as follows:
\begin{itemize}
\item a positive residue at a fixed coupling is \LE{} evidence that the atom is present at that
  coupling, not evidence for the continuum gap;
\item the dissolution of the signal at strong coupling is \emph{consistent} with the analytic
  picture (the state is at or beyond the cutoff), not evidence against P3;
\item a one-sided interval that includes zero is the honest summary of a probe whose object is
  not exactly the gauge-invariant $0^{++}$;
\item a diagnostic that separates atom from cut on synthetic data is a method, not a proof.
\end{itemize}
Every number below carries one of these readings.

\subsection{The $\mathrm{SU}(2)$ residue where it is resolvable}

At $\beta=2.2$ on $\mathrm{SU}(2)$, $L=6$, $N=4000$ configurations, the connected zero-momentum
$\operatorname{tr}F^2$ correlator gives
\[
f_{1,\mathrm{bare}}^2=7.14\pm0.42\quad(\text{seventeen sigma}),\qquad
m_G=1.51\pm0.15\,a ,
\]
with every smearing level strongly positive. This is \LE. It is the clearest direct sighting of
the residue in the campaign. At $\beta=1.0$ the signal is $0.2\sigma$ and the residue is not
extractable. The reason is not a failure of the residue but a property of the strong-coupling
lattice: $m\cdot a\gtrsim3$, so on $T=6$ the state sits at or beyond the cutoff. P3 is an
analytic statement in the regime where $m\cdot a$ diverges, so the state is provably present but
numerically sub-lattice-spacing. Strong-coupling provability and numerical positivity live at
opposite ends of the coupling; the two observations are consistent.

\subsection{The $k\to0$ curvature probe, read honestly}

On a single $L=8$ volume, the $\mathrm{SU}(3)$ $\beta=5.7$ softest nonzero mode
$k=2\pi/L=\pi/4$ gives
\[
G''=+0.307\pm0.055\quad(\text{6/6 transverse plane-wave modes}),\qquad
W_{\mathrm{geo}}=+0.319\pm0.005 ,
\]
with the action bracket at equipartition ($+0.037\pm0.099$). $\mathrm{SU}(2)$ at $\beta=2.4$
reproduces the sign ($+0.192\pm0.058$, $W_{\mathrm{geo}}=+0.252$). The measured points are
positive and volume-stable, and $W_{\mathrm{geo}}$ is pinned at $+0.32$.

\textbf{The caveat that must not be dropped.} This is a $k\to0$ sign probe on a
\emph{gauge-dependent gluon mode} in a flat chart, not the gauge-invariant $0^{++}$ residue.
Keeping only the proven $W_{\mathrm{geo}}\ge0$ and treating the rest statistically gives $95\%$
one-sided lower bounds of $-0.098$ ($\mathrm{SU}(3)$) and $-0.035$ ($\mathrm{SU}(2)$):
\textbf{zero is not excluded.} Closing it is one variance inequality, $(\blacklozenge)=$ H1, not
a better dataset. The gauge-invariant object is
$m_{\mathrm{GI}}=\inf\operatorname{supp}\rho_s$ of $s=(\beta/2g)\operatorname{tr}F^2$.

\subsection{The volume confound is removed}

At fixed $k=\pi/2$, $L=4$ and $L=8$ agree ($+0.425\pm0.084$ versus $+0.396\pm0.063$); at fixed
$L=8$, $G''(\pi/4)=+0.307<G''(\pi/2)=+0.396$ but does not collapse. The earlier $L=4\to6$ drop
mixed a $k$ decrease with a volume change; it is now separated. $W_{\mathrm{geo}}$ is pinned at
$+0.32$ across $L=4,6,8$ and both $k$. This is a genuine control: the sign of the probe is not a
volume artifact.

\figprobe

\subsection{The right-object audit}

A post-audit measured the true gauge-invariant $0^{++}$ object --- the $k=0$ mode of
$q(x)=1-\langle\tfrac12\operatorname{Re}\operatorname{tr}U_p\rangle\approx\operatorname{tr}F^2$
--- rather than the gauge-dependent gluon ($\mathrm{SU}(2)$, $L=8$, $N=200$). Two findings, both
negative for the variance route:
\begin{enumerate}
\item The composite zero-momentum variance scales as $\operatorname{Var}(\bar q)\propto V^{-1.0}$
  ($L=4,6,8$). So the variance form of H1, $\operatorname{Var}\le1/\lambda$, is vacuously
  satisfied by any local composite at large $V$ and carries no gap content; its apparent
  $G''>0$ is the wrong scaling (a susceptibility $1/\operatorname{Var}\propto V$, not an inverse
  propagator).
\item The gauge-dependent gluon mode, which does have the right variance behaviour, fails the
  free-field acid test at every $\beta$ measured ($G''/\text{free}=0.18$--$0.63$; one-sided lower
  bound $<0$ at every $\beta$).
\end{enumerate}
Consequently the $D/S\approx2\%$ quoted in earlier drafts is a compactness/Haar number measured
on that gauge-dependent mode --- not gap evidence --- and is not robust under a susceptibility
normalisation ($-0.45$ to $+0.60$ at strong coupling; the $2\%$ is a weak-coupling value). The
physical datum is the glueball two-point mass, which is \emph{not} resolved at $L=8$ (fit
$m_0\gtrsim\pi$, $\xi\lesssim a$). The audit sharpens, and does not close, the gap; it is
consistent with $G''_a=a^2m^2$ and with H1 being necessary, not sufficient.

\subsection{The Savvidy lattice cell}

Chapter~\ref{sec:savvidy} gives the structural reading; here are the numbers. The defect ratio
$D/S$ is bounded ($\le1.1$) but $O(1)$: $0.94\le D/S\le1.00$. Five of twelve modes have $G''<0$.
The $k\to0$ mode is absent on the available volumes. The earlier reported $G''\approx+0.28$ is a
small-volume artifact of the free dispersion $4\sin^2(\pi/L)$. The lattice does not contradict
the theorem-set, and it does not confirm a positive margin. It shows that the margin is thin.

\subsection{The diagnostics}

The campaign's strongest methodological output on the numerical side is a set of diagnostics for
the atom that are genuine and continuum-robust \emph{in the ideal}, and gated by $L\to\infty$ in
practice. They are worth stating precisely because they define what a decisive lattice
computation would look like.

\subsubsection{The P3 approach law}

Define
\begin{equation}
T(\tau)=-\frac{d}{d\tau}\log\bigl|m(\tau)-\hat e_1(\tau)\bigr| ,
\label{eq:approachlaw}
\end{equation}
where $m(\tau)$ is a two-exponential fit to the correlator at Euclidean time $\tau$ and
$\hat e_1$ its leading effective mass. For an isolated atom, $T(\tau)\to\Delta$ (the isolation
gap); for a pure continuum, $T(\tau)\to1/\tau$. The statement is \LT{} within the stated model
(Watson's lemma: an atom below a threshold gives an exponential rate, a pure continuum gives a
power law). It is robust to a background continuum and to the choice of operator
(\code{WI-P3-CONTINUUM.md}, \code{WI-P3-RIGOR.md}). \textbf{But}: on a single finite volume every
discrete spectrum is a sum of exponentials, so a single $L$ cannot decide; the discriminating
content is the $L\to\infty$ scaling of the approach, which is isolation $=$ H1.

\subsubsection{The one-sided isolation bound}

The quantitative version is a one-sided bound
$\Delta\ge\kappa(\tau)\to\Delta$, sharper than the qualitative approach law
(\code{WI-P3-QUANT.md}). It is \LT{} given the data, but it needs the mass $e_1$ as an input and
is window-limited: a thin $\Delta$ is under-certified. It reads, it does not force.

\subsubsection{The Re-slope statistic}

For exact complex-$\tau$ zeros of the correlator's analytic continuation, the statistic
\begin{equation}
R=\frac{d\operatorname{Re}\tau}{d\ln c_1}=\frac{1}{E_1}
\label{eq:reslope}
\end{equation}
is robust to a background continuum: an atom on a gapped or threshold continuum exposes five to
six atom-slope rungs, a pure cut none (\code{WI-MODIFIED-FUNCTIONAL.md}, \code{WI-RESLOPE-LATTICE.md}).
It is the campaign's clearest methodological gain on the analytic line. \textbf{But} exact
complex-$\tau$ zeros are A3 (a reading), and on finite real-$\tau$ data the statistic is blind at
$O(h)$ and collapses to the P3 plateau (\code{WI-RESLOPE-FINITE.md}). So it is a read, gated by
A3 and by more data than is currently available.

\subsubsection{The finite-size law}

The scaling
\begin{equation}
\Delta(L)=m_1+\frac{c}{m_1L^2}
\qquad\text{for an atom, versus}\qquad
\Delta(L)\sim\frac{(2\pi/L)^2}{m}\to0\quad\text{for a cut}
\label{eq:fss}
\end{equation}
gives a uniform criterion: $\Delta(L,a)\ge\delta$ for all $L$ iff $m_1>0$, which is the gap
(\code{WI-P3-UNIFORM.md}). It is the cleanest statement of what a definitive lattice computation
must demonstrate, and it is also the cleanest statement of why a single volume is insufficient.

\begin{table}[t]
\centering
\footnotesize
\begin{tabularx}{\textwidth}{@{}>{\raggedright\arraybackslash}p{3.2cm}
  >{\raggedright\arraybackslash}p{4.1cm}>{\raggedright\arraybackslash}X@{}}
\toprule
\textbf{Diagnostic} & \textbf{What it measures} & \textbf{Status / limiting step}\\
\midrule
P3 approach law $T(\tau)$ & atom: $T\to\Delta$; cut: $T\to1/\tau$ & \LT{} in the model;
  robust to continuum and operator; single $L$ cannot decide; $L\to\infty=\,$H1\\
Isolation bound $\Delta\ge\kappa(\tau)$ & a one-sided quantitative gap & \LT{} given data;
  needs the edge; window-limited\\
Re-slope $R=1/E_1$ & atom-slope rungs along $\operatorname{Im}\tau$ & \LT{} with exact
  complex-$\tau$ zeros; A3 in practice; blind on real-$\tau$ data\\
Finite-size law $\Delta(L)=m_1+c/(m_1L^2)$ & the $L\to\infty$ limit & \LT{} criterion;
  uniformity is $m_1>0=$ the gap\\
Two-exponential fit & $f_1^2$, $m_1$ at fixed $L$ & \LE; finite-volume fact; not the
  continuum atom\\
\bottomrule
\end{tabularx}
\caption{The diagnostics and their exact logical status.}
\label{tab:diag}
\end{table}

\figdiag

\subsection{Statistical methods}

The numbers in this chapter are only as good as their error analysis, and the error analysis is
part of the honest reading. The protocol is standard, but it is worth stating because the
difference between a detection and a bound is exactly the difference between two statistical
statements.

\textbf{Autocorrelation and binning.} Configurations generated by a local update are correlated.
The integrated autocorrelation time $\tau_{\mathrm{int}}$ of the observable is estimated from the
normalized autocorrelation function (summed until it first crosses zero, with the standard
window), and configurations are grouped into bins of length $\gg\tau_{\mathrm{int}}$ before any
average is formed. A residual correlation shows up as a bin-size dependence of the error; the
protocol requires the error to be stable under doubling the bin size.

\textbf{Jackknife.} All errors quoted are jackknife errors over bins. For a ratio or a fit
parameter (such as the residue from a two-exponential fit), the jackknife is performed at the
level of the fit, not by propagating errors through the fit.

\textbf{Detection versus bound.} A detection is a statement that a central value is many
standard deviations from zero: $f_{1,\mathrm{bare}}^2=7.14\pm0.42$ is a $17\sigma$ detection at
that coupling and volume. A bound is a statement that a one-sided confidence interval excludes
zero: the $k\to0$ probe gives a $95\%$ one-sided lower bound of $-0.098$ (SU(3)), which does
\emph{not} exclude zero. The two statements answer different questions, and the paper never
converts one into the other. In particular, a detection on a gauge-dependent mode is not a bound
on the gauge-invariant object.

\textbf{Systematics.} The dominant systematic of the $k\to0$ probe is the object choice (gluon
mode versus composite); the dominant systematic of the residue measurement is the fit range and
the smearing family. Where a systematic cannot be quantified, the paper reports the spread across
choices rather than a single number, and the number is labelled \LE{} with the caveat.

\textbf{Finite-volume discreteness.} On a finite lattice the spectrum is discrete, so every
correlator is a finite sum of exponentials and a two-exponential fit will always ``work'' over
some range. This is why the diagnostics of \S\ref{ch:numerics} insist on a volume-scaling law
rather than a single-volume fit: the question is not whether an exponential fits, but whether its
rate survives $L\to\infty$.

\subsection{What a decisive numerical test would need}

The diagnostics make the requirements explicit. A decisive lattice computation would need:
\begin{enumerate}
\item several volumes at fixed coupling, large enough that the $1/L^2$ law is visible and the
  lowest momentum is in the scaling regime (the present $L=8$ is not);
\item a gauge-invariant $0^{++}$ operator whose two-point function is measured with a smearing
  family, so that the residue and the mass are operator-independent;
\item a continuum extrapolation at fixed physical volume, with at least three lattice spacings;
\item a one-sided statistical statement (a confidence bound) that excludes zero for
  $\Delta(L,a)$ uniformly, not merely a central value;
\item a systematic-error budget for the continuum extrapolation and for the scale setting.
\end{enumerate}
The campaign's throughput analysis (Appendix~\ref{app:protocols}) indicates that $\mathrm{SU}(3)$
$L=16$--$24$ at three couplings is feasible in C/OpenMP on 16 cores over days, and faster on a
GPU; the bottleneck is not the algorithm but the statistical precision needed for a one-sided
bound at the percent level. The existing $\mathrm{SU}(2)$ and $\mathrm{SU}(3)$ data are
consistent with the picture but do not meet these requirements.

\subsection{Summary of the numerical status}

The numerics do three things, and no more. They \emph{see} the residue where it is resolvable
($17\sigma$ at $\beta=2.2$ on $\mathrm{SU}(2)$). They \emph{locate} the sign of the $k\to0$
curvature on a gauge-dependent mode (positive, volume-stable, but with a one-sided interval that
includes zero). They \emph{calibrate} the diagnostics (P3, Re-slope, FSS) on synthetic data. They
do not prove the gap, and they are not offered as evidence for it. Their value is that they turn
each open step into a falsifiable numerical question, which is the most a finite-volume
computation can do for a continuum problem.
